\documentclass[11pt]{article}

\usepackage{amsmath, amsthm, amssymb}
\usepackage{geometry}
\geometry{margin=1in}

\newtheorem{theorem}{Theorem}
\newtheorem{corollary}{Corollary}

\title{Metric Collapse as the Structural Driver of Entropy Depth}
\author{Inacio Vasquez}
\date{}

\begin{document}

\maketitle

\begin{abstract}
We classify entropy-depth lower bounds for local refinement processes
via chromatic properties of power graphs.
We prove that bounded-degree fixed-radius systems admit constant depth,
while linear entropy depth arises precisely from metric collapse
of $2R$-neighborhood geometry.
\end{abstract}

\section{Setup}

Let $G_n$ be a connected graph.
For locality radius $R$, define the power graph $G_n^{(2R)}$
with edges between vertices at distance $\le 2R$.

Depth satisfies:
\[
T_n \ge \chi(G_n^{(2R)}).
\]

\section{Bounded Local Regime}

\begin{theorem}
If $\Delta(G_n)=O(1)$ and $R=O(1)$,
then $\chi(G_n^{(2R)}) = O(1)$.
\end{theorem}

\begin{proof}
The $2R$-ball size is bounded by
\[
1 + \Delta \sum_{i=0}^{2R-1} (\Delta-1)^i.
\]
Thus maximum degree in $G_n^{(2R)}$ is bounded.
Greedy coloring completes the proof.
\end{proof}

\section{Metric Collapse Regime}

\begin{theorem}
If $\operatorname{diam}(G_n) \le 2R$,
then $G_n^{(2R)} = K_n$ and
$\chi(G_n^{(2R)}) = n$.
\end{theorem}

\section{Logarithmic Collapse in Small-Diameter Families}

\begin{theorem}
If $\operatorname{diam}(G_n)=O(\log n)$
and $R=\Theta(\log n)$,
then $\chi(G_n^{(2R)})=\Theta(n)$.
\end{theorem}

\section{Unified Dichotomy}

Exactly one holds:

\begin{itemize}
\item Fixed $R$, bounded degree $\Rightarrow T_n = O(1)$.
\item $R \ge \tfrac{1}{2}\operatorname{diam}(G_n)$ $\Rightarrow T_n = n$.
\item $R=\Theta(\log n)$ in expanders $\Rightarrow T_n=\Theta(n)$.
\end{itemize}

\section{Conclusion}

Entropy depth is governed by $2R$-metric collapse.
Expansion alone does not force linear depth.
Linear entropy depth arises precisely when locality radius
compresses global metric structure.

\end{document}
